% UBT-AI-PROVENANCE-BEGIN
% schema: ubt-ai-provenance/v1
% tier: B_machine_verified
% ai_assistance: disclosed
% human_review: machine-verification
% editorial_responsibility: Ing. David Jaroš
% policy: ../../AI_PROVENANCE.md
% notice: Machine-verified against named sources or verifiers; individual attestation is not claimed.
% UBT-AI-PROVENANCE-END

\ifdefined\INCLUDEMODE\else
\documentclass[11pt]{article}
\usepackage[a4paper,margin=2.5cm]{geometry}
\usepackage{fontspec}
\usepackage{amsmath,amssymb}
\usepackage[hidelinks]{hyperref}
\usepackage{xurl}
\emergencystretch=3em
\usepackage{ubtprovenance}
\UBTTier{B}
\renewcommand{\UBTProvenanceNotice}{\par\smallskip{\footnotesize Provenience AI --- úroveň B. Text vznikl s pomocí AI; strojově ověřen uvedeným ověřovačem; individuální lidské potvrzení netvrdíme.}\par\smallskip}
\title{Krok 6: opravený stav předpokladů QED}
\author{UBT}
\date{2026-10-31}
\begin{document}\maketitle\UBTProvenanceNotice
\fi

% BILINGUAL-UNIT: qed-summary.status
\section{Oprava závislostí}
GAP-U1-NEUTRALITY: FORMER PROOF REFUTED; PHYSICAL BRIDGE OPEN. Původní důkaz neutrality pro levé působení fáze je neplatný. Bezhmotnost na nenulovém pozadí není dokázána. Skalární srovnávací model místo toho vyžaduje následující podmínku; nejde o odvozenou kinetickou akci UBT.
\[m_\gamma^2=2e^2\lVert Qv\rVert^2,\qquad m_\gamma^2=0\Longleftrightarrow Qv=0\quad(e\ne0).\]

% BILINGUAL-UNIT: qed-summary.downstream
\section{Navazující výsledky}
Kroky 2--5 zůstávají nepárovými staršími výpočty s nevyřešenými předpoklady, nikoli ověřenými odvozeními na nabitém nenulovém pozadí UBT. Standardní srovnání s jedno-smyčkovým Schwingerovým členem QED je následující výraz. Jeho numerické vyčíslení neodvozuje z UBT bezhmotný foton, Diracův vrchol, neutrální skalár, Yukawovu vazbu ani vakuum. Dřívější výrazy s chybným faktorem dvě a paušální dvousmyčkové odhady ve starším kroku 4 nejsou autoritativní. Yukawova interakce neutrálního skaláru sama může přispět již v jedné smyčce; neutralita neodstraňuje všechny skalární smyčkové efekty. Bez hmotnosti a vazeb z ní neplyne univerzální numerická korekce.
\[a_e^{(1),\mathrm{QED}}=\frac{\alpha}{2\pi}.\]

Odkaz pro smyčku neutrálního skaláru: \url{https://arxiv.org/abs/1403.2309}.

% BILINGUAL-UNIT: qed-summary.checks
\section{Ověření a otevřená podmínka}
Opravený krok 1 a párová výzkumná poznámka uvádějí předpoklady, přesný protipříklad, normalizaci hmotnosti, odvození odezvy složené veličiny a záznam strojového ověření. Numerická pomůcka \path{tools/verify_qed_phi_const.py} je pouze podmíněným srovnáním. Vyžaduje výslovný náboj a nemůže potvrdit uzavření UBT. LEAN-PENDING: není nainstalován lean ani lake. Výchozím jazykem překladu je angličtina; netvrdíme významové schválení ani lidské potvrzení.
\ifdefined\INCLUDEMODE\def\UBTFinish{}\else\def\UBTFinish{\end{document}}\fi
\UBTFinish
